% ECONOMICS_PROBLEM: {"id": "I11-346", "title": "Social value and pricing of pharmaceutical innovation", "jel_code": "I11", "source_id": "OP-0051"}
% !TeX program = xelatex
\documentclass[11pt,a4paper]{article}
\usepackage[margin=23mm]{geometry}
\usepackage{fontspec}
\setmainfont{Latin Modern Roman}
\usepackage{amsmath,amssymb,mathtools}
\usepackage{xcolor,enumitem,booktabs,longtable,array,xurl,hyperref}
\definecolor{muted}{HTML}{53636C}
\hypersetup{colorlinks=true,urlcolor=blue,linkcolor=blue}
\setlength{\parindent}{0pt}
\setlength{\parskip}{5pt}
\newcommand{\R}{\mathbb R}
\newcommand{\E}{\mathbb E}
\newcommand{\N}{\mathbb N}
\newcommand{\ind}{\mathbf 1}
\newcommand{\Var}{\operatorname{Var}}
\newcommand{\Cov}{\operatorname{Cov}}
\newcommand{\supp}{\operatorname{supp}}
\DeclareMathOperator*{\argmin}{arg\,min}
\DeclareMathOperator*{\argmax}{arg\,max}
\newcommand{\entrymeta}[3]{{\sffamily\small\color{muted}\textbf{\texttt{#1}}\quad|\quad #2\par #3\par}}
\newcommand{\Problemref}[1]{\texttt{#1}}
\newcommand{\JELref}[2]{\texttt{#2} (JEL \texttt{#1})}
\begin{document}
\section*{Social value and pricing of pharmaceutical innovation}
\textbf{Problem ID:} \texttt{I11-346}\quad \textbf{JEL:} \texttt{I11}\par
% BEGIN ECONOMICS STATEMENT
\label{problem:OP-0051}
\label{atlas:prob:H1}
\noindent\textbf{Original atlas identifier:} H1; entry 51. \textbf{Field:} Health economics. \textbf{Original tractability label:} Hard.

\noindent\textbf{Classification:} Parameterized theoretical/empirical research agenda; not a certified unsolved theorem.

\subsection*{Definitions and assumptions}
Fix dates $t=0,\ldots,T$, indications $k=1,\ldots,K$, and countries $j=1,\ldots,J$. A pricing/reward policy $p$ specifies prices, reimbursement, patent duration or buyouts, and any subsidies. A dynamic model $m$ determines firms' indication-specific research $r_{kt}$, innovation arrival probabilities, product quality, entry, and access $x_{jkt}$. Let $Q_{jt}(p;m)$ be quality-adjusted life-years gained relative to a stated baseline, $v_j\geq0$ their monetary social value, and $L_{jt}(p;m)$ other real health/economic benefits or harms, valued consistently. Let $C_{jt}$ denote real treatment/production costs and $C_{Rt}$ research resource costs. Fix discount factor $\beta\in(0,1]$ and nonnegative country weights $w_j$. Specify supplier and payer equilibrium selection and account for all countries whose R\&D or access responds.

\subsection*{Formal research question}
\emph{Editorial formulation: the mathematical target below is not a theorem quoted from the references.}

Characterize the feasible-policy welfare frontier
\[
 W_m(p)=\mathbb E_m\sum_{t=0}^{T}\beta^t\left\{
 \sum_{j=1}^{J}w_j[v_jQ_{jt}+L_{jt}-C_{jt}]-C_{Rt}\right\}.
\]
All quantities on the right depend on $p$ and $m$. Prices are transfers, not additional resource costs; distributional gains/losses from transfers require explicit payer/producer welfare weights. Impose payer budgets, participation, commitment, safety/approval constraints, and a funding rule. Identify how indication-level revenue shocks affect research allocation and eventually clinical health gains, rather than only patent or approval counts. Compare price caps, reimbursement rules, and buyouts using $\varepsilon$-optimal policies over an identified model set, retaining uncertainty about innovation timing and quality.

\subsection*{Interpretation and scope}
Static patient access and future innovation are distinct margins, and the same price reduction can improve the first while changing the second in either direction under portfolio substitution. Market-size shocks may affect scientific opportunity, disease prevalence, or competing indications; an exclusion restriction is needed before treating them as exogenous revenue variation. A buyout must specify valuation, licensing, and collusion controls rather than presume a known patent value. Country weights and the value of health are normative inputs; global innovation spillovers cannot be summarized by domestic consumer surplus. Costs should include failed research projects and displaced scientific capacity. The substantive extension connects causal revenue responses to eventual global health value and feasible contracting, while existing innovation-market-failure and patent-buyout models remain solved benchmarks.

\subsection*{Source audit and status}
[S1]--[S3] are identifiable primary works: general innovation theory, a buyout mechanism, and pharmaceutical market-size evidence. They do not uniquely establish the atlas's broad question as an open theorem or verify its status in 2026. The objective above is an editorial social-accounting formulation.

\subsection*{References}
\begin{enumerate}[label={[S\arabic*]},leftmargin=*,itemsep=0.6em]
\item Kenneth J. Arrow (1962). Economic Welfare and the Allocation of Resources for Invention. In The Rate and Direction of Inventive Activity: Economic and Social Factors, 609--626. Princeton University Press for NBER.
\url{https://www.nber.org/books-and-chapters/rate-and-direction-inventive-activity-economic-and-social-factors/economic-welfare-and-allocation-resources-invention}
\newline\emph{Locator and support limits:} Chapter: innovation-market failures; no pharmaceutical-specific dynamic price optimum.
\item Michael Kremer (1998). Patent Buyouts: A Mechanism for Encouraging Innovation. The Quarterly Journal of Economics 113(4), 1137--1167. DOI: 10.1162/003355398555865.
\url{https://academic.oup.com/qje/article-abstract/113/4/1137/1916997}
\newline\emph{Locator and support limits:} Abstract and mechanism: patent-buyout design; effects depend on information and collusion assumptions.
\item Daron Acemoglu and Joshua Linn (2004). Market Size in Innovation: Theory and Evidence from the Pharmaceutical Industry. The Quarterly Journal of Economics 119(3), 1049--1090. DOI: 10.1162/0033553041502144.
\url{https://academic.oup.com/qje/article-abstract/119/3/1049/1938820}
\newline\emph{Locator and support limits:} Abstract and article: market-size and pharmaceutical innovation; product counts do not by themselves measure health surplus.
\end{enumerate}

\subsection*{Common conventions: Source mathematical conventions}

Unless an entry states otherwise, agent and firm sets are finite. State and action spaces are measurable, strategies and outcome maps are measurable, probability laws are specified, and expectations exist for the targets under discussion. Infinite-horizon discount factors lie in $(0,1)$ and discounted objectives are finite. Norms, tolerances and complexity input sizes must be specified for computational comparisons. Constants or primitives that appear locally are inputs, not empirical facts asserted by this catalogue.

\textbf{Identification.} A statistical model $m\in\mathcal M$ implies an observable law $P_m$ and target $\theta(m)$. At law $P$, its identified set is
\[
\Theta(P;\mathcal M)=\{\theta(m):m\in\mathcal M,\ P_m=P\}.
\]
Point identification holds when this set is a singleton, or equivalently when $P_{m_1}=P_{m_2}$ implies $\theta(m_1)=\theta(m_2)$ for all compatible models. Sharp bounds mean bounds attained, or approached when the boundary is not attained, by compatible models. Writing this set is a definition, not a solution: the substantive task is to establish informative restrictions and derive or estimate the set. An unrestricted-confounding model may give only trivial bounds.

\textbf{Inference.} Identification, estimability and valid uncertainty quantification are separate questions. Fix $\alpha\in(0,1)$. A confidence procedure $C_n$ over a declared law class $\mathcal P$ has asymptotic uniform coverage at level $1-\alpha$ when
\[
\liminf_{n\to\infty}\inf_{P\in\mathcal P}\Pr_P\{\theta(P)\in C_n\}\geq1-\alpha.
\]
For set-identified targets the coverage event must be defined separately, for example coverage of the full identified set. If an entry uses identification-indexed models instead of uniquely indexed laws, the same coverage convention is applied over those models. Pointwise limits do not establish uniform validity, and asymptotic validity does not guarantee finite-sample precision. Sample sizes, dependence structures and weak-identification sequences are part of each application.

\textbf{Counterfactuals.} $Y_i(a)$ is a potential outcome under a specified intervention, including its timing, implementation and versions. It is not $Y_i$ conditional on voluntarily choosing $a$. A full intervention vector permits interference; shorthand scalar treatment notation does not silently impose no interference. Assignment, selection, attrition, anticipation and measurement must be specified. A claim about an instrument's local effect is not automatically a population policy effect.

\textbf{Economic equilibrium.} A counterfactual model includes best responses, constraints, feasibility, market clearing, state transitions and expectations consistent with the stated information. When several equilibria exist, either a declared selector is an input or the target is set-valued. Comparisons across policy regimes must use the stated selection convention. Reduced-form relationships are not presumed invariant after an equilibrium-changing intervention.

\textbf{Optimization and welfare.} Optimization is over the stated feasible set. If attainment is not established, $\sup$ or $\inf$ and $\varepsilon$-optimal choices replace an asserted maximizer or minimizer. For example, with value $V=\sup_{a\in\mathcal A}W(a)<\infty$, an $\varepsilon$-optimal set is $\{a\in\mathcal A:W(a)\geq V-\varepsilon\}$ for $\varepsilon>0$. Benefits, costs, risk and health or environmental outcomes can be aggregated only using declared units and conversion weights. Interpersonal utility weights, the treatment of fiscal transfers and foreign profits, discounting, and the behavioral welfare criterion are explicit inputs. A comparative-static derivative additionally requires existence and appropriate differentiability of the selected counterfactual.

\textbf{Sources.} Local labels [S1], [S2], and so forth restart in every question. Locators identify the relevant abstract, model, section or result at the level actually checked; a bibliographic or abstract-level verification is not represented as inspection of an entire proof. Working papers are distinguished from later publications and changing versions. Source summaries are paraphrased. All URL checks and editorial formulations refer to the audit date stated above.
% END ECONOMICS STATEMENT
\end{document}
